% ============================================================
%  cap04.tex  --  Capítulo 4: Séries Infinitas (Perpetuidades)
%  Baseado em: Notas de Aula 3 (Aula 4) + Notas de Aula 4 (complemento)
%  Referência: Bueno, Rangel & Santos, Cap. 4
% ============================================================

\chapter{Séries Infinitas -- Perpetuidades}

\section{Motivação}

\begin{definicao}
Uma \textbf{perpetuidade} é uma série de pagamentos que se estende indefinidamente no tempo
($n \to \infty$). O problema central é: quanto vale \emph{hoje} uma renda que dura para sempre?
\end{definicao}

Perpetuidades surgem na avaliação de ações (Modelo de Gordon), fundos de endowment, rendas
vitalícias e na análise de dívidas soberanas (consols britânicos).

% ─────────────────────────────────────────────────────────────────────────────
\section{Séries Discretas}

\subsection{Perpetuidade Uniforme}

Partindo do valor presente de uma série vencida com $n$ períodos:
\[
  P = R\left[\frac{(1+i)^n-1}{i}\right]\frac{1}{(1+i)^n}
    = \frac{R}{i}\left[1 - \frac{1}{(1+i)^n}\right].
\]
Tomando o limite $n \to \infty$ (e como $(1+i)^n \to \infty$ para $i>0$):

\begin{resultado}
\[
  \boxed{P = \frac{R}{i}.}
  \qquad\text{Equivalentemente: } R = P\cdot i.
\]
A lógica é direta: os juros $P \cdot i$ gerados pelo capital $P$ a cada período são exatamente
retirados sem jamais tocar no principal.
\end{resultado}

\begin{exemplo}{Renda perpétua de R\$\,1.500/mês}
\label{ex:perpetuidade-uniforme}
Com taxa de $0{,}5\%$ a.m., o capital necessário para gerar R\$\,1.500/mês para sempre:
\[
  P = \frac{1{.}500}{0{,}005} = \textbf{R\$\,300.000.}
\]
\end{exemplo}

\subsection{Perpetuidade em PA (Gradiente Crescente)}

Quando os pagamentos crescem em PA -- $\{R, 2R, 3R, \ldots\}$ -- o valor presente finito é:
\[
  P_n = \frac{R}{i}\left\{(1+i)\left[\frac{(1+i)^n-1}{i}\right]\cdot\frac{1}{(1+i)^n}
        - \frac{n}{(1+i)^n}\right\}.
\]
Usando a Regra de L'Hôpital para calcular os limites quando $n\to\infty$
(ambas as parcelas $\to (1+i)/i$ e $0$ respectivamente):

\begin{resultado}
\[
  P = \frac{R\,(1+i)}{i^2}.
\]
\end{resultado}

\begin{exemplo}{Perpetuidade em PA}
$R = 100$, $\Delta = 100$ (gradiente), $i = 2\%$ a.m.:
\[
  P = \frac{100 \times 1{,}02}{(0{,}02)^2} = \frac{102}{0{,}0004}
    = \textbf{R\$\,255.000.}
\]
Interpretação: pagar hoje R\$\,255.000 equivale a receber R\$\,100 no mês~1, R\$\,200 no mês~2, e assim por diante, sem fim.
\end{exemplo}

\begin{nota}
\textbf{Nota 4.1 -- Séries decrescentes não convergem.}
Para uma série em PA \emph{decrescente} os pagamentos eventualmente se tornam negativos;
a série não possui sentido econômico como perpetuidade.
Para séries em PG decrescentes ($\alpha < 0$), a convergência requer $|\alpha| < i$ (ver Seção~\ref{sec:pgperp}).
\end{nota}

\subsection{Perpetuidade em PG}
\label{sec:pgperp}

Para pagamentos em PG com crescimento $\alpha$ por período:
\[
  P_n = R\left[\frac{1-\left(\frac{1+\alpha}{1+i}\right)^n}{i-\alpha}\right].
\]
Para convergência é necessário $\alpha < i$. Tomando $n \to \infty$:

\begin{resultado}
\[
  \boxed{P = \frac{R}{i - \alpha}.}
\]
\end{resultado}

\begin{obs}
Este resultado é válido para $i > \alpha$. Se $\alpha \geq i$, a série diverge
(pagamentos crescem mais rápido do que a taxa de desconto) e o valor presente é infinito.
\end{obs}

% ─────────────────────────────────────────────────────────────────────────────
\section{Séries Contínuas}

\subsection{Perpetuidade Uniforme Contínua}

Para um fluxo contínuo $b$ por unidade de tempo a uma taxa contínua $r$:
\[
  P = \lim_{T\to\infty} \frac{b(1-e^{-rT})}{r} = \frac{b}{r}.
\]

\begin{resultado}
\[
  P = \frac{b}{r}.
\]
\end{resultado}

\begin{exemplo}{Comparação discreta vs.\ contínua (Exemplo 4.7)}
$b = R = 1{.}500$ por mês, $r = \ln(1{,}005) \approx 0{,}4988\%$ a.m.:
\[
  P_{\text{cont}} = \frac{1{.}500}{0{,}004988} \approx \textbf{R\$\,300.720.}
\]
Praticamente idêntico ao resultado discreto (R\$\,300.000), pois o intervalo é pequeno.

\smallskip
\textit{Motivo:} O resultado contínuo expressa o somatório de taxas infinitesimais que não proporcionam aquela retirada (período = taxa).
\end{exemplo}

\subsection{Perpetuidade Gradiente Crescente Contínua}

Para fluxo $b(t) = b_0 \, e^{\alpha t}$ (crescimento contínuo a taxa $\alpha$):
\[
  P = \int_0^\infty b_0 e^{\alpha t} e^{-rt}\,dt
    = b_0 \int_0^\infty e^{-(r-\alpha)t}\,dt
    = \frac{b_0}{r - \alpha}, \qquad r > \alpha.
\]

\begin{resultado}
\[
  P = \frac{b_0}{r - \alpha}.
\]
\end{resultado}

\subsection{Perpetuidade Geométrica Discreta Finita (revisão)}

Para $n$ períodos:
\[
  P_n = R\left[\frac{1 - e^{(\alpha-r)n}}{r - \alpha}\right].
\]
Quando $n\to\infty$ e $r>\alpha$: $P = R/(r-\alpha)$.

\begin{exemplo}{Exemplo 4.9 -- perpetuidade geométrica contínua}
$R = 100$, $r_{\text{am}} = 1{,}0\%$ a.m., $\alpha = 0{,}5\%$ a.m.:
\[
  P = \frac{100}{0{,}010 - 0{,}005} = \frac{100}{0{,}005} = \textbf{R\$\,20.000.}
\]
\end{exemplo}

% ─────────────────────────────────────────────────────────────────────────────
\section{Modelo de Gordon}

O Modelo de Gordon (1956) é a aplicação direta da fórmula de perpetuidade geométrica à avaliação
de ações.

\begin{definicao}
O \textbf{preço de uma ação} é o valor presente do fluxo de dividendos futuros esperados,
descontados pela taxa de retorno requerida (custo de oportunidade) $i$:
\[
  P = \frac{D}{(1+i)} + \frac{D}{(1+i)^2} + \cdots
\]
\end{definicao}

\paragraph{Caso uniforme (dividendo constante $D$):}
\[
  P = \frac{D}{i}.
\]

\paragraph{Caso geométrico (dividendo cresce a $\alpha$ por período, $i > \alpha$):}
\[
  P_0 = \frac{D_0(1+\alpha)}{(1+i)} + \frac{D_0(1+\alpha)^2}{(1+i)^2} + \cdots
\]

\begin{resultado}
\[
  \boxed{P_0 = \frac{D_0(1+\alpha)}{i - \alpha} = \frac{D_1}{i - \alpha},}
\]
onde $D_1 = D_0(1+\alpha)$ é o dividendo esperado no próximo período.
\end{resultado}

\begin{moderno}[title={Modelo de Gordon -- Petrobras (PETR4)}]
Hipóteses ilustrativas (dados de referência, 2024):
\begin{itemize}[leftmargin=1.8em, nosep]
  \item Dividendo pago no último exercício: $D_0 = \text{R\$\,}3{,}20$ por ação.
  \item Taxa de crescimento estimada dos dividendos: $\alpha = 5\%$ a.a.
  \item Taxa de retorno requerida pelo acionista: $i = 15\%$ a.a. (CAPM).
\end{itemize}
\[
  P_0 = \frac{3{,}20 \times 1{,}05}{0{,}15 - 0{,}05}
       = \frac{3{,}36}{0{,}10} = \textbf{R\$\,33{,}60.}
\]
Se o preço de mercado for inferior a R\$\,33,60, o ativo está subavaliado pelo modelo.
Note a sensibilidade: se $\alpha$ subir para $7\%$, o preço-justo sobe para R\$\,42,24 ($+25\%$).
\end{moderno}
